\subsection{拓扑扭转}
\label{sec:twist}

理论不仅可以表述在（半）Minkowski 空间上。例如，我们可以考虑圆柱或带形，仍取 Minkowski 度规，
它们给出闭弦或开弦的态。我们也可以通过时间线的 Wick 转动，把这些流形上的理论表述为欧氏度规下的理论。
进一步，我们可以通过标准的协变化，把系统表述在带有任意度规与自旋结构的二维流形 $\Sigma$ 上。
% [原文 p.5]
超对称还能存活吗？我们当然可以通过对 $Q_\pm\cO$ 与 $\overline Q_\pm\cO$ 的协变化、并把变分参数
$\epsilon_\pm$ 与 $\bar\epsilon_\pm$ 取为自旋丛 $S_\pm$ 的截面，来推广场的超对称变换的定义
\begin{equation}
  \delta\cO=i\epsilon_+Q_-\cO-i\epsilon_-Q_+\cO-i\bar\epsilon_+\overline Q_-\cO+i\bar\epsilon_-\overline Q_+\cO. \tag{2.12}
\end{equation}
然而，协变化后作用量的不变性要求变分参数协变常值，这在弯曲流形上是不可能的。有几种办法可以恢复
一部分超对称。其中之一就是我们现在要描述的\emph{拓扑扭转}。

设 $F=F_V$ 或 $F_A$ 守恒，并具有\emph{荷整性}，即它生成一个 $U(1)$ 对称群，使得非旋量场与旋量场
分别具有偶的与奇的荷。当我们考虑定向 Riemann 流形 $\Sigma$ 上的理论时，拓扑扭转就是把一个
R-荷为 $q$、取值于向量丛 $E$ 的场，替换为取值于 $E\otimes T_\Sigma^{\otimes q/2}$ 的场。
这里 $T_\Sigma$ 是装配了 Levi-Civita 联络的全纯切丛。注意，由于 R-荷宇称的限制，
$E\otimes T_\Sigma^{\otimes q/2}$ 无需选择 $\Sigma$ 的自旋结构便有定义。变分参数 $\epsilon_\pm$
与 $\bar\epsilon_\pm$ 也发生同样的改变，其中一些变成标量。我们可以把这样的参数取为常数，
于是相应的超荷守恒。对 $F=F_V$ 称为 \emph{A-扭转}，对 $F=F_A$ 称为 \emph{B-扭转}。
在 A-扭转（相应地 B-扭转）理论中，超荷 $\overline Q_+$ 与 $Q_-$、从而它们的和 $Q=Q_A$
（相应地 $\overline Q_\pm$ 及其和 $Q=Q_B$）守恒。此外，$Q$-闭算子的关联函数只依赖于算子的
$Q$-上同调类，并且在 $\Sigma$ 的 Riemann 度规形变下不变。特别地，它们只依赖于 $\Sigma$ 的拓扑
——我们得到一个拓扑场论。

我们还可以把扭转后的理论与某个称为\emph{拓扑引力}的二维引力理论耦合，从而考虑\emph{拓扑弦论}。
$g$-圈振幅由对亏格 $g$ 曲线的模空间的积分给出。

\subsection{例子}
\label{sec:examples}

% [Tutor Note: 拓扑扭转的关键是 R-荷的整性条件：只有 e^{πiF} 是某个规范群元素时，
% E⊗T_Σ^{⊗q/2} 才能在有向 Riemann 曲面上良定义而不需要自旋结构。A-扭转（F=F_V）保留
% Q=Q_A=\overline Q_++Q_-，B-扭转（F=F_A）保留 Q=Q_B=\overline Q_++\overline Q_-；
% 这与 §2.2 中 A-膜/B-膜所保留的超荷一致，也正是 §2.7 镜像自同构 Q_-↔\overline Q_-、
% F_V↔F_A 的来源。]

\para{非线性 sigma 模型}

设 $X=(X,g)$ 是紧 Kähler 流形。则存在一个以 $X$ 为靶的二维 $(2,2)$ 超对称 QFT，称为
\emph{非线性 sigma 模型}。作为数据的一部分，我们还选取一个类 $[B]\in H^2(X,\R/2\pi\Z)$，
称为 \emph{B-场}。

% [原文 p.6]
曲面 $\Sigma$ 上的场变量由一个映射 $\phi:\Sigma\to X$ 和一个取值于 $\phi^*TX$ 的 Dirac 费米子
$\psi$ 组成。在 Minkowski 时空（$\Sigma=\R^2$，时间与空间坐标为 $x^0$ 与 $x^1$）上，
拉氏量为\footnote{在本讲义中，我们采用作用量为拉氏量密度的积分除以 $2\pi$ 的约定：
\begin{equation}
  S=\frac{1}{2\pi}\int_\Sigma \dd^2x\,\cL. \tag{2.13}
\end{equation}
B-场以 $-i\int_\Sigma \phi^*B$ 进入欧氏作用量，等价地，它以 $\exp\big(i\int_\Sigma\phi^*B\big)$
进入路径积分权重。这解释了 B-场的周期性。}
\begin{align}
  \cL &= g_{i\bar j}\Big(\partial_0\phi^i\partial_0\bar\phi^{\bar j}-\partial_1\phi^i\partial_1\bar\phi^{\bar j}
   +i\bar\psi^{\bar j}_-(D_0+D_1)\psi^i_-+i\bar\psi^{\bar j}_+(D_0-D_1)\psi^i_+\Big)\notag\\
  &\quad +R_{i\bar j k\bar l}\psi^i_+\psi^k_-\bar\psi^{\bar j}_-\bar\psi^{\bar l}_+ + 2\pi(\phi^*B)_{10}. \tag{2.14}
\end{align}
这里我们选取了 $X$ 的局部全纯坐标 $z^i$，而 $\phi^i$ 等是分量 $\phi^i=z^i\circ\phi$ 等等。
$D_\mu\psi^i_\pm=\partial_\mu\psi^i_\pm+\partial_\mu\phi^j\Gamma^i_{jk}\psi^k_\pm$ 是关于 $TX$ 上
Levi-Civita 联络的拉回而言的协变导数，而
$R_{i\bar jk\bar l}=-g_{\bar l m}(R_{i\bar j})^{m}{}_{k}$ 是 $TX$ 的 Riemann 曲率。
拉氏量在超对称变换
$\delta=i\epsilon_+Q_- - i\epsilon_-Q_+ - i\bar\epsilon_+\overline Q_- + i\bar\epsilon_-\overline Q_+$ 下不变：
\begin{align}
  \delta\phi^i &= \epsilon_+\psi^i_--\epsilon_-\psi^i_+, \tag{2.15}\\
  \delta\psi^i_\pm &= \pm i\bar\epsilon_\mp(\partial_0\pm\partial_1)\phi^i+\epsilon_\pm\Gamma^i_{jk}\psi^j_+\psi^k_-. \tag{2.16}
\end{align}
复共轭变量的变换由这些变换按规则 $(AB)^\dagger=B^\dagger A^\dagger$ 得到。例如，
$\delta\bar\phi^{\bar\imath}=\bar\psi^{\bar\imath}_-\bar\epsilon_+-\bar\psi^{\bar\imath}_+\bar\epsilon_-=-\bar\epsilon_+\bar\psi^{\bar\imath}_-+\bar\epsilon_-\bar\psi^{\bar\imath}_+$。

该模型经典地同时具有矢量与轴 $U(1)$ R-对称性且具荷整性，但若第一 Chern 类 $c_1(X)$ 非零，
则轴 R-对称性有异常：轴转动把 $[B]$ 平移 $c_1(X)$。该模型经典地是标度不变的，但靶度规在 RG 流下
改变。Kähler 类\footnote{我们按 $\omega=\frac{i}{2\pi}g_{i\bar j}\dd z^i\wedge\dd\bar z^{\bar j}$
归一化 Kähler 类，使得对全纯映射 $\phi:\Sigma\to X$，路径积分权重为
$\exp\big(-\int_\Sigma\phi^*(\omega-iB)\big)$。}%
按 Ricci 流跑动：当能标从 $\mu$ 变到 $\mu'$ 时，
$[\omega]\to[\omega']=[\omega]+c_1(X)\log(\mu'/\mu)$。当 $c_1(X)=0$ 时，理论在红外极限下流向
一个中心荷为 $\hat c=\dim_\C X$ 的 SCFT。该模型的手征环等为
\[
\begin{array}{rcl@{\qquad}l}
 \cR_A &=& \mathrm{QH}^*(X) & \text{量子上同调环},\\
 \cR_B &=& \mathrm{H}^*(X,\wedge^*T_X) & \text{多向量场的上同调环},\\
 \cC_A &=& \mathrm{Fuk}(X) & \text{Fukaya 范畴},\\
 \cC_B &=& \mathrm{D}^b_{\mathrm{Coh}}(X) & \text{具凝聚上同调的层的导出范畴}.
\end{array}
\]
